%% ma: L10-B14-0056
%% lop: 10
%% chuong: 5
%% bai: 14
%% dang: 10.14.12
%% loai: TL
%% muc: VD
%% dap_an: "-"
%% nguon: "chuyen-de-cac-so-dac-trung-cua-mau-so-lieu-khong-ghep-nhom-toan-10-chuong-trinh-moi.pdf (GV. Phan Nhật Linh) (P05), Bài 03 Dạng 2 Bài tập tự luận 16"
%% kien_thuc: "Lớp 10 Bài 13 (số trung bình, trung vị, tứ phân vị, mốt), Bài 14 (khoảng biến thiên, khoảng tứ phân vị, phương sai, độ lệch chuẩn, giá trị bất thường)"
%% trang_thai: chinh-thuc
%% ghi_chu: "Lời giải gốc sai phương sai nhà máy B (6,55; đúng ≈ 5,80) và kết luận nhà máy A lương cao hơn (sai vì số trung bình bị kéo bởi giá trị bất thường); ý c) gốc lặp lại ý a), đã bỏ. Tự duyệt theo yêu cầu giáo viên 22/10/2026."
\begin{ex}
Kết quả điều tra mức lương hằng tháng (đơn vị: triệu đồng) của một số công nhân ở hai nhà máy A và B như sau:
\begin{center}
\begin{tabular}{|c|c|c|c|c|c|c|c|c|c|}\hline
Nhà máy A & $4$ & $5$ & $5$ & $47$ & $5$ & $6$ & $4$ & $4$ & \\\hline
Nhà máy B & $2$ & $9$ & $9$ & $8$ & $10$ & $9$ & $9$ & $11$ & $9$\\\hline
\end{tabular}
\end{center}
a) Hãy tìm số trung bình, mốt, các tứ phân vị và độ lệch chuẩn của hai mẫu số liệu.

b) Hãy tìm các giá trị bất thường trong mỗi mẫu số liệu. Công nhân nhà máy nào có mức lương cao hơn? Tại sao?
\loigiai{a) Nhà máy A: sắp xếp $4$; $4$; $4$; $5$; $5$; $5$; $6$; $47$. $\overline x_A=10$; mốt $4$ và $5$; $Q_1=4$, $Q_2=5$, $Q_3=5{,}5$; $s_A^2=\dfrac{1}{8}\sum x_i^2-10^2=196$, $s_A=14$.

Nhà máy B: sắp xếp $2$; $8$; $9$; $9$; $9$; $9$; $9$; $10$; $11$. $\overline x_B=\dfrac{76}9\approx8{,}44$; mốt $9$; $Q_2=9$, $Q_1=\dfrac{8+9}2=8{,}5$, $Q_3=\dfrac{9+10}2=9{,}5$; $s_B^2=\dfrac{694}9-\left(\dfrac{76}9\right)^2\approx5{,}80$, $s_B\approx2{,}41$.

b) A: $\Delta_Q=1{,}5$; $Q_3+1{,}5\Delta_Q=7{,}75$; $Q_1-1{,}5\Delta_Q=1{,}75$ nên $47$ là giá trị bất thường.

B: $\Delta_Q=1$; $Q_3+1{,}5\Delta_Q=11$; $Q_1-1{,}5\Delta_Q=7$ nên $2$ là giá trị bất thường.

Số trung bình của A lớn hơn chỉ do giá trị bất thường $47$. So sánh theo trung vị ($5$ và $9$) hoặc mốt thì công nhân nhà máy B có mức lương cao hơn.}
\end{ex}
